% Document/LinearAlgebra/VectorSpaces/LinearMaps/HomSpaces.tex
% Section: Hom Spaces as Vector Spaces

\newpage
\section{Hom Spaces as Vector Spaces}

\begin{theorem}[{$\Hom[K]{U}{V}$ is a Vector Space}]
\label{thm:hom-vs}
    For $K$-vector spaces $U$ and $V$, the set $\Hom[K]{U}{V}$ of all $K$-linear maps
    $T: U \to V$ is a $K$-vector space under pointwise operations:
    \begin{align*}
        (S + T)(u) &= S(u) + T(u) \\
        (aT)(u) &= a \cdot T(u)
    \end{align*}
\end{theorem}

\begin{proof}
    We verify all vector space axioms.
    
    \textbf{Closure under addition}:
    For $S, T \in \Hom[K]{U}{V}$, we show $S + T \in \Hom[K]{U}{V}$:
    \begin{align*}
        (S + T)(au_1 + bu_2) &= S(au_1 + bu_2) + T(au_1 + bu_2) \\
        &= aS(u_1) + bS(u_2) + aT(u_1) + bT(u_2) \\
        &= a(S(u_1) + T(u_1)) + b(S(u_2) + T(u_2)) \\
        &= a(S+T)(u_1) + b(S+T)(u_2)
    \end{align*}
    
    \textbf{Closure under scalar multiplication}:
    For $a \in K$ and $T \in \Hom[K]{U}{V}$, we show $aT \in \Hom[K]{U}{V}$:
    \begin{align*}
        (aT)(bu_1 + cu_2) &= a \cdot T(bu_1 + cu_2) \\
        &= a \cdot (bT(u_1) + cT(u_2)) \\
        &= abT(u_1) + acT(u_2) \\
        &= b(aT)(u_1) + c(aT)(u_2)
    \end{align*}
    
    \textbf{Additive identity}:
    The zero map $0: U \to V$ defined by $0(u) = 0_V$ is in $\Hom[K]{U}{V}$, and
    $(T + 0)(u) = T(u) + 0_V = T(u)$, so $T + 0 = T$.
    
    \textbf{Additive inverse}:
    For $T \in \Hom[K]{U}{V}$, define $(-T)(u) = -T(u)$. Then $-T \in \Hom[K]{U}{V}$ since
    $(-T)(au + bv) = -(aT(u) + bT(v)) = a(-T)(u) + b(-T)(v)$, and
    $(T + (-T))(u) = T(u) + (-T(u)) = 0_V$, so $T + (-T) = 0$.
    
    \textbf{Commutativity of addition}:
    $(S + T)(u) = S(u) + T(u) = T(u) + S(u) = (T + S)(u)$, using commutativity in $V$.
    
    \textbf{Associativity of addition}:
    $((R + S) + T)(u) = (R+S)(u) + T(u) = R(u) + S(u) + T(u) = R(u) + (S+T)(u) = (R + (S+T))(u)$.
    
    \textbf{Compatibility of scalar and field multiplication}:
    $((ab)T)(u) = (ab)T(u) = a(bT(u)) = a((bT)(u)) = (a(bT))(u)$, so $(ab)T = a(bT)$.
    
    \textbf{Identity element of scalar multiplication}:
    $(1_K T)(u) = 1_K \cdot T(u) = T(u)$, so $1_K T = T$.
    
    \textbf{Distributivity of scalar multiplication over vector addition}:
    $(a(S + T))(u) = a(S+T)(u) = a(S(u) + T(u)) = aS(u) + aT(u) = (aS + aT)(u)$.
    
    \textbf{Distributivity of scalar multiplication over field addition}:
    $((a + b)T)(u) = (a+b)T(u) = aT(u) + bT(u) = (aT)(u) + (bT)(u) = (aT + bT)(u)$.
\end{proof}

\newpage
\section{Endomorphism Algebras}

\begin{definition}[Endomorphism Space]
\label{def:end}
    $\End[K]{V} = \Hom[K]{V}{V}$ is the space of all linear maps from $V$ to itself.
    It carries an additional multiplication structure: composition of maps.
    That is, for $S, T \in \End[K]{V}$, we define $S \cdot T = \Composition{S}{T}$.
\end{definition}

\begin{definition}[$K$-Algebra]
\label{def:k-algebra}
    A \textbf{$K$-algebra} is a $K$-vector space $A$ equipped with a bilinear multiplication
    $\cdot: A \times A \to A$ such that $(A, +, \cdot)$ is a ring and the following
    \textbf{compatibility} holds:
    \[
        a(xy) = (ax)y = x(ay) \quad \Forall :{a \in K, x, y \in A}{}
    \]
    
    A $K$-algebra is \textbf{associative} if $(xy)z = x(yz)$, and \textbf{unital} if 
    it has a multiplicative identity $1_A$.
\end{definition}

\begin{theorem}[{$\End[K]{V}$ is a Ring}]
\label{thm:end-ring}
    $\End[K]{V}$ is a ring with composition as multiplication:
    \begin{itemize}
        \item Addition: $(S + T)(v) = S(v) + T(v)$
        \item Multiplication: $(S \cdot T)(v) = S(T(v)) = (\Composition{S}{T})(v)$
        \item Zero: $0(v) = 0_V$
        \item Unity: $\Identity{V}(v) = v$
    \end{itemize}
\end{theorem}

\begin{proof}
    $(\End[K]{V}, +)$ is abelian (Theorem~\ref{thm:hom-vs}).
    
    \textbf{Associativity}: $(\Composition{R}{S}) \cdot T = \Composition{R}{(\Composition{S}{T})}$ (function composition).
    
    \textbf{Distributivity}: $(\Composition{R}{(S + T)})(v) = R(S(v) + T(v)) = R(S(v)) + R(T(v))$.
    
    \textbf{Unity}: $\Composition{\Identity{V}}{T} = \Composition{T}{\Identity{V}} = T$.
\end{proof}

\begin{theorem}[{$\End[K]{V}$ is a $K$-Algebra}]
\label{thm:end-algebra}
    $\End[K]{V}$ is a unital associative $K$-algebra.
\end{theorem}

\begin{proof}
    Vector space and ring structures are established. For compatibility:
    \[
        ((\Composition{aS}{T}))(v) = (aS)(T(v)) = a \cdot S(T(v)) = a \cdot (\Composition{S}{T})(v)
    \]
    Similarly $(\Composition{S}{(aT)})(v) = S(a \cdot T(v)) = a \cdot S(T(v))$.
\end{proof}

\begin{theorem}[{Non-Commutativity of $\End[K]{V}$}]
\label{thm:end-noncommutative}
    If $\Dim[K]{V} \geq 2$, then $\End[K]{V}$ is non-commutative.
\end{theorem}

\begin{proof}
    Since $\Dim[K]{V} \geq 2$, there exist linearly independent vectors $v_0, v_1 \in V$.
    Extend $\{v_0, v_1\}$ to a basis $B$ of $V$ (by Zorn's lemma if $V$ is infinite-dimensional).
    
    Define $S, T \in \End[K]{V}$ via their action on $B$:
    \[
        S(v_0) = v_1, \quad S(v_1) = 0, \quad S(b) = 0 \text{ for } b \in B \setminus \{v_0, v_1\}
    \]
    \[
        T(v_0) = 0, \quad T(v_1) = v_0, \quad T(b) = 0 \text{ for } b \in B \setminus \{v_0, v_1\}
    \]
    Extend linearly to all of $V$.
    
    \textbf{Compute $\Composition{S}{T}$}:
    $(\Composition{S}{T})(v_0) = S(T(v_0)) = S(0) = 0$.
    $(\Composition{S}{T})(v_1) = S(T(v_1)) = S(v_0) = v_1$.
    
    \textbf{Compute $\Composition{T}{S}$}:
    $(\Composition{T}{S})(v_0) = T(S(v_0)) = T(v_1) = v_0$.
    $(\Composition{T}{S})(v_1) = T(S(v_1)) = T(0) = 0$.
    
    Thus $(\Composition{S}{T})(v_0) = 0 \neq v_0 = (\Composition{T}{S})(v_0)$, so $\Composition{S}{T} \neq \Composition{T}{S}$.
\end{proof}

\newpage
\section{Dimension of Hom Spaces}

\begin{theorem}[{Dimension of $\Hom[K]{U}{V}$ (Finite-Dimensional Case)}]
\label{thm:hom-dim-finite}
    Let $U$ and $V$ be finite-dimensional $K$-vector spaces with
    $\Dim[K]{U} = m$ and $\Dim[K]{V} = n$. Then
    \[
        \Dim[K]{\Hom[K]{U}{V}} = m \cdot n
    \]
\end{theorem}

\begin{proof}
    Let $(u_0, \ldots, u_{m-1})$ be a basis of $U$ and $(v_0, \ldots, v_{n-1})$ be a basis of $V$.
    For each pair $(i, j)$ with $0 \leq i \leq m-1$ and $0 \leq j \leq n-1$, define
    $T_{ij} \in \Hom[K]{U}{V}$ by its values on the basis:
    \[
        T_{ij}(u_k) = \delta_{ik} \, v_j =
        \begin{cases}
            v_j & \text{if } k = i \\
            0   & \text{if } k \neq i
        \end{cases}
    \]
    and extend linearly. Each $T_{ij}$ is well-defined and linear by the universal property
    of bases. We claim $\{T_{ij} \mid 0 \leq i \leq m-1,\; 0 \leq j \leq n-1\}$ is a basis
    for $\Hom[K]{U}{V}$.

    \textbf{Spanning.}
    Let $T \in \Hom[K]{U}{V}$. Since $(v_0, \ldots, v_{n-1})$ is a basis of $V$, for each
    $i \in \{0, \ldots, m-1\}$ there exist unique scalars $a_{ij} \in K$ such that
    \[
        T(u_i) = \sum_{j \in n} a_{ij} \, v_j
    \]
    Define $S = \sum_{i \in m} \sum_{j \in n} a_{ij} \, T_{ij}$. Then for each basis
    vector $u_k$:
    \[
        S(u_k) = \sum_{i \in m} \sum_{j \in n} a_{ij} \, T_{ij}(u_k)
        = \sum_{i \in m} \sum_{j \in n} a_{ij} \, \delta_{ik} \, v_j
        = \sum_{j \in n} a_{kj} \, v_j
        = T(u_k)
    \]
    Since $S$ and $T$ agree on a basis of $U$, by the uniqueness of the extension, $S = T$.
    Hence $T = \sum_{i,j} a_{ij} T_{ij}$ is a linear combination of the $T_{ij}$'s.

    \textbf{Linear independence.}
    Suppose $\sum_{i \in m} \sum_{j \in n} a_{ij} \, T_{ij} = 0$ (the zero map).
    Evaluating at $u_k$ for each $k \in \{0, \ldots, m-1\}$:
    \[
        0 = \sum_{i \in m} \sum_{j \in n} a_{ij} \, T_{ij}(u_k)
        = \sum_{j \in n} a_{kj} \, v_j
    \]
    Since $(v_0, \ldots, v_{n-1})$ is linearly independent, $a_{kj} = 0$ for all
    $j \in \{0, \ldots, n-1\}$. As this holds for every $k$, all coefficients $a_{ij} = 0$.

    Thus $\{T_{ij}\}$ is a basis of $\Hom[K]{U}{V}$ with $m \cdot n$ elements.
\end{proof}

\begin{corollary}[Dimension of End]\label{cor:end-dim-early}
    If $\Dim[K]{V} = n$, then $\Dim[K]{\End[K]{V}} = n^2$.
\end{corollary}

\begin{proof}
    $\Dim[K]{\End[K]{V}} = \Dim[K]{\Hom[K]{V}{V}} = n \cdot n = n^2$.
\end{proof}
